\chapter{子流形几何与活动标架}
\label{ch:13-submanifold-geometry}
\ChapterPlan{周4.5（本地页码待核对）；梁上册第二版4.4及文内练习；Andrews第17--19讲补充}{正文及证明初稿；8题详解与3幅原创图，尚未独立数学审阅}

同一个黎曼流形可以有不同的等距浸入。度量刻画内在长度，第二基本形式则记录切向量在环境中求导时离开切空间的部分。本章把第2章曲面的分解推广到任意余维，推导三组相容方程，再用活动标架反过来构造局部欧氏浸入。最后通过体积第一变分解释平均曲率向量与极小性的关系。

设$f:M^m\to(\bar M^{m+r},\bar g)$为光滑浸入，$m\geq1$，环境和诱导度量均正定、光滑，除变分区域外不考虑边界。始终取$g=f^*\bar g$，并把$TM$经$\dd f$识别为$f^*T\bar M$的子丛。不同原像点即使有相同像点，也保留各自的切空间和法空间。横线表示环境对象；曲率仍取$R(X,Y)=\nabla_X\nabla_Y-\nabla_Y\nabla_X-\nabla_{[X,Y]}$。本章不把Lorentz超曲面的法向符号直接代入正定公式。

\Needspace{17\baselineskip}
\section{Gauss与Weingarten公式}
\subsection*{沿浸入求导，而非在自交像上选延拓}
环境联络在$f^*T\bar M$上诱导沿$f$的联络，记$\bar\nabla_Xs$。若$y^A$为环境坐标、$s=s^A\partial_{y^A}|_f$，则
\[
 (\bar\nabla_Xs)^A=X(s^A)+\bar\Gamma^A_{BC}(f)X(f^B)s^C.
\]
第10章的换基规则保证其与坐标无关。它可以沿每条曲线由环境协变导数定义，故不要求$s$在整个自交像上有单值延拓。其曲率为$\bar R(\dd fX,\dd fY)s$，以下简记$\bar R(X,Y)s$。

\begin{definition}[法丛与第二基本形式]\label{def:normal-second-form}
法丛$NM\subset f^*T\bar M$的纤维为$N_pM=(\dd f(T_pM))^\perp$。用$\top,\perp$表示正交投影，定义
\[
 \nabla_XY=(\bar\nabla_XY)^\top,\qquad
 B(X,Y)=(\bar\nabla_XY)^\perp,\qquad
 \nabla_X^\perp\xi=(\bar\nabla_X\xi)^\perp.
\]
$B$称法丛值第二基本形式；$\nabla^\perp$称法联络。对法向场$\xi$定义形算子$A_\xi X=-(\bar\nabla_X\xi)^\top$。
\end{definition}

局部取切标架$u_i$，其Gram矩阵$g_{ij}$可逆，投影为$P^\top(s)=g^{ij}\bar g(s,u_i)u_j$，因而光滑。$P^\perp=I-P^\top$有常秩$r$；第6章常秩子丛结论给出$NM$的光滑结构，且
\[
 f^*T\bar M=TM\mathbin{\oplus^\perp}NM.
\]
法丛不一定有全局单位标架；本章只在小邻域选取局部标架。

\begin{theorem}[正交分解与诱导联络]\label{thm:gauss-weingarten}
$\nabla$为$g$的Levi--Civita联络，$\nabla^\perp$与法丛内积相容，$B\in\Gamma(\operatorname{Sym}^2T^*M\otimes NM)$。Gauss和Weingarten公式为
\[
 \bar\nabla_XY=\nabla_XY+B(X,Y),\qquad
 \bar\nabla_X\xi=-A_\xi X+\nabla_X^\perp\xi,
\]
其中$g(A_\xi X,Y)=\bar g(B(X,Y),\xi)$。$A$对$X,\xi$均为函数线性，且每个$A_\xi$自伴。
\end{theorem}
\begin{proof}
切向投影保留联络的方向线性和Leibniz规则。由环境度量相容，
\[
 Xg(Y,Z)=g(\nabla_XY,Z)+g(Y,\nabla_XZ),
\]
因为法向项与切向量正交。环境无挠给
$\bar\nabla_X(\dd fY)-\bar\nabla_Y(\dd fX)=\dd f[X,Y]$：坐标式中混合偏导抵消，$\bar\Gamma$下槽对称。切投影说明$\nabla$无挠，由第11章唯一性即为Levi--Civita；法投影说明$B(X,Y)=B(Y,X)$。

对函数$a,b$展开得$\bar\nabla_{aX}(bY)=ab\bar\nabla_XY+aX(b)Y$，最后一项法投影为零，所以$B(aX,bY)=abB(X,Y)$。对法向场同理，$\bar\nabla_X(a\xi)=X(a)\xi+a\bar\nabla_X\xi$中的首项切投影为零，故$A$张量性成立。对$\bar g(\xi,Y)=0$求导，
\[
 0=\bar g(\bar\nabla_X\xi,Y)+\bar g(\xi,\bar\nabla_XY)
   =-g(A_\xi X,Y)+\bar g(\xi,B(X,Y)).
\]
$B$对称给$A_\xi$自伴。最后对两法向场内积求导并投影，得到$\nabla^\perp$的度量相容性。
\end{proof}

\subsection*{曲线加速度、全测地与超曲面}
沿$M$中曲线$c$，分解公式为
\[
 \bar D_t(\dd f\dot c)=\dd f(D_t\dot c)+B(\dot c,\dot c).
\]
在坐标中逐项对$\dot c^i\dd f(\partial_i)$求导即可得到，包括变速项。因此$M$的测地线在环境中一般有法向加速度。
\begin{figure}[htbp]
\centering\DGDiagram{submanifold-acceleration}

\caption{加速度在切空间与法空间中的正交分解。图用平面曲线$F(u)=(u,u^2/2)$及$u(t)=t+t^2/2$在$t=0$的精确数据：环境加速度为$(1,1)$，切向、法向分量分别为$(1,0)$、$(0,1)$；该参数曲线不是内蕴测地线。}
\label{fig:submanifold-acceleration}
\end{figure}

\begin{proposition}[全测地的局部判据]\label{prop:totally-geodesic}
$B=0$当且仅当$M$的每条充分短测地线经$f$成为具有同一仿射参数的环境测地线。这时称$f$全测地。此结论不要求完备性，也不保证全局距离相等。
\end{proposition}
\begin{proof}
正向由加速度分解立即得到。反向，对任意$p,v$取$M$的短测地线初值$(p,v)$，环境加速度为零给$B_p(v,v)=0$。由对称性和极化
$2B(u,v)=B(u+v,u+v)-B(u,u)-B(v,v)$知$B=0$。这只比较局部微分方程，不声称其像在环境中实现任意两端点的距离下确界。
\end{proof}

当$r=1$，局部选单位法向$n$，则$\nabla^\perp_Xn=0$，因为它既在一维法空间中，又由$\bar g(n,n)=1$得与$n$正交。令$II(X,Y)=\bar g(B(X,Y),n)$，便有
\[
 B=II\,n,\quad A_n=S,\quad \bar\nabla_Xn=-SX,
 \quad II(X,Y)=g(SX,Y).
\]
欧氏环境下恰为第2章的$S=-\dd n$。换$n$为$-n$时$II,S$变号，但$B$不变。在高余维，不同$A_\xi$未必交换，不能默认存在使它们同时对角化的主方向基。

\begin{definition}[平均曲率向量]\label{def:mean-curvature-vector}
对任意切空间正交归一基$(e_i)_{i=1}^m$，定义
\[
 \mathbf H=\frac1m\operatorname{tr}_g B
 =\frac1m\sum_{i=1}^m B(e_i,e_i).
\]
正交换基消去交叉项，故与基无关。$\mathbf H=0$称极小浸入；$B=0$则为更强的全测地条件。
\end{definition}
超曲面上$\mathbf H=Hn$，$H=\operatorname{tr}S/m$；当$m=2$与第2章完全一致。球面$S^m(a)\subset\R^{m+1}$取外法向$n=f/a$时$S=-I/a$、$B=-gn/a$、$\mathbf H=-n/a$，向量指向球心。维数因子$m$将在体积变分中显式出现。


\Needspace{17\baselineskip}
\section{Gauss--Codazzi--Ricci方程}
给$B$求导时，值域用法联络，两个输入槽用切联络。具体定义
\[
 (\nabla B)(X;Y,Z)=\nabla_X^\perp(B(Y,Z))
 -B(\nabla_XY,Z)-B(Y,\nabla_XZ).
\]
这是$T^*M\otimes\operatorname{Sym}^2T^*M\otimes NM$截面；若只对$B$的分量作普通偏导，会丢失三组联络项。法曲率取
$R^\perp(X,Y)\xi=\nabla_X^\perp\nabla_Y^\perp\xi-\nabla_Y^\perp\nabla_X^\perp\xi-\nabla_{[X,Y]}^\perp\xi$。

\Needspace{20\baselineskip}
\begin{theorem}[三组相容方程]\label{thm:gauss-codazzi-ricci}
对切向量场$X,Y,Z,W$和法向场$\xi,\eta$，有
\begin{align*}
 g(R(X,Y)Z,W)&=\bar g(\bar R(X,Y)Z,W)\\
 &\quad+\bar g(B(X,W),B(Y,Z))-\bar g(B(Y,W),B(X,Z)),\\
 (\bar R(X,Y)Z)^\perp&=(\nabla B)(X;Y,Z)-(\nabla B)(Y;X,Z),\\
 \bar g(R^\perp(X,Y)\xi,\eta)&=\bar g(\bar R(X,Y)\xi,\eta)
       +g([A_\xi,A_\eta]X,Y).
\end{align*}
其中$[A_\xi,A_\eta]=A_\xi A_\eta-A_\eta A_\xi$。三式依次为Gauss、Codazzi、Ricci方程；末式不是Ricci张量的缩并公式。
\end{theorem}
\begin{proof}
在$\bar\nabla_X\bar\nabla_YZ$中先代入Gauss公式，再分别对切、法两部分求导，得到
\begin{align*}
 (\bar\nabla_X\bar\nabla_YZ)^\top
   &=\nabla_X\nabla_YZ-A_{B(Y,Z)}X,\\
 (\bar\nabla_X\bar\nabla_YZ)^\perp
   &=B(X,\nabla_YZ)+\nabla_X^\perp(B(Y,Z)).
\end{align*}
减去$X,Y$交换项及$\bar\nabla_{[X,Y]}Z$。切向结果为
\[
 (\bar R(X,Y)Z)^\top=R(X,Y)Z-A_{B(Y,Z)}X+A_{B(X,Z)}Y.
\]
移项并与$W$取内积即得Gauss式。法向结果减去两份$\nabla B$后，剩下
$B(\nabla_XY-\nabla_YX-[X,Y],Z)=0$，故得Codazzi式。

再对$\xi$使用Weingarten公式，两次求导的法向部分为
\[
 (\bar\nabla_X\bar\nabla_Y\xi)^\perp
 =-B(X,A_\xi Y)+\nabla_X^\perp\nabla_Y^\perp\xi.
\]
反对称化给
\[
 (\bar R(X,Y)\xi)^\perp
 =R^\perp(X,Y)\xi-B(X,A_\xi Y)+B(Y,A_\xi X).
\]
后两项与$\eta$配对为
$-g(A_\eta X,A_\xi Y)+g(A_\eta Y,A_\xi X)
=-g([A_\xi,A_\eta]X,Y)$，整理即得第三式。整个推导在任意标架成立，不要求预先选择主方向。
\end{proof}

\begin{corollary}[截面曲率与欧氏缩并]\label{cor:submanifold-curvature-traces}
若$u,v$正交归一，则
\[
 K_M(u,v)=\bar K(u,v)+\bar g(B(u,u),B(v,v))-|B(u,v)|^2.
\]
特别对欧氏浸入，令$|B|^2=\sum_{i,j}|B(e_i,e_j)|^2$，有
\[
 \Ric(X,X)=m\bar g(B(X,X),\mathbf H)-\sum_i|B(X,e_i)|^2,
 \qquad\mathrm{Scal}=m^2|\mathbf H|^2-|B|^2.
\]
因此欧氏极小子流形的$\Ric\leq0$、$\mathrm{Scal}=-|B|^2\leq0$；不由此断言其每个截面曲率均非正。
\end{corollary}
\begin{proof}
在Gauss式代入$(X,Y,Z,W)=(u,v,v,u)$即得首式。欧氏环境曲率为零，再将$(e_i,X,X,e_i)$代入并对$i$求和，第一项成为$\bar g(B(X,X),\sum_iB(e_i,e_i))$，第二项为所列平方和。再取$X=e_j$并求和，得到标量式。极小时去掉含$\mathbf H$项即可。维数大于2时，若干截面曲率的和非正，不足以逐项确定符号。
\end{proof}

\subsection*{超曲面、法平坦与环境曲率项}
在常截面曲率$c$的环境中，$\bar R(X,Y)Z=c(\bar g(Y,Z)X-\bar g(X,Z)Y)$。对超曲面，在一点取$S$的正交特征基，$Se_i=k_i e_i$，于是
\[
 K_M(e_i,e_j)=c+k_i k_j\quad(i\ne j),\qquad
 (\nabla_XS)Y=(\nabla_YS)X.
\]
后一式来自环境曲率的法向部分为零及$\nabla^\perp n=0$。一般环境中不能略去这一法向项。正球面外法向给$k_i=-1/a$，仍有$K=a^{-2}$；圆柱面一主曲率为零，故内蕴平坦。

在欧氏环境中，Ricci方程说明
\[
 R^\perp=0\quad\Longleftrightarrow\quad
 [A_\xi,A_\eta]=0\text{ 对所有法向量 }\xi,\eta.
\]
因为右式与所有$X,Y$配对为零等价于算子为零。有限组可交换自伴算子可同时正交对角化：先分解一个算子的特征空间，其余算子因交换保持各特征空间，再在这些空间内递归对角化。取法空间有限基即可应用。因此法平坦给逐点共同主方向；它不自动给一个处处光滑、甚至全局定义的主方向标架。超曲面的法联络局部单位标架中恒零，所以法曲率恒零；但全局法丛的平行绕行仍可能变号。

\Needspace{16\baselineskip}
\begin{proposition}[欧氏全脐超曲面的局部分类]\label{prop:umbilic-classification}
设$m\geq2$，欧氏超曲面在连通邻域选有单位法向$n$且$S=\lambda I$。则$\lambda$为常数；$\lambda=0$时像局部为仿射超平面的开集，$\lambda\ne0$时为半径$1/|\lambda|$球面的开集。未假设完备时不能将开集改为整个球面或平面。
\end{proposition}
\begin{proof}
Codazzi式给$(X\lambda)Y=(Y\lambda)X$。在任意点，给单位$X$选与之正交的单位$Y$，与$Y$取内积得$X\lambda=0$，故$\dd\lambda=0$。连通性给常数。

若$\lambda=0$，则$\dd n=0$，$n=n_0$恒定，且$\dd\langle f,n_0\rangle=0$，故像在某仿射超平面中。若$\lambda\ne0$，则
\[
 \dd(f+\lambda^{-1}n)=\dd f+\lambda^{-1}\dd n=0,
\]
所以$C=f+\lambda^{-1}n$恒定，$|f-C|=1/|\lambda|$。两种情形中，$f$到相应$m$维超曲面的微分均为同维同构，逆映射定理说明它在每点附近是到开集的微分同胚。维数1时$S=\lambda I$对每条曲线自动成立，却不迫使曲率恒定。
\end{proof}


\Needspace{17\baselineskip}
\section{活动标架与Cartan结构方程}
\subsection*{适配正交标架的分块联络}
局部取切丛正交归一标架$e_i$和法丛正交归一标架$n_\alpha$，$1\leq i,j\leq m$、$1\leq\alpha,\beta\leq r$。记$\theta^i$为切余标架，定义一形式矩阵
\[
 \omega^i{}_j(X)=g(\nabla_Xe_j,e_i),\quad
 \tau^\alpha{}_\beta(X)=\bar g(\nabla_X^\perp n_\beta,n_\alpha),\quad
 \beta^\alpha{}_i(X)=\bar g(B(X,e_i),n_\alpha).
\]
这里$\beta$是矩阵名；其分量的上标$\alpha$是法向指标。于是
\[
 \bar\omega=\begin{pmatrix}\omega&-\beta^{\mathsf T}\\
                   \beta&\tau\end{pmatrix},\qquad
 \omega^{\mathsf T}=-\omega,\quad\tau^{\mathsf T}=-\tau,
 \quad\beta^\alpha{}_i=h^\alpha_{ij}\theta^j,
 \quad h^\alpha_{ij}=h^\alpha_{ji}.
\]
联络矩阵的列指标始终表示被求导的基向量，例如$\bar\nabla e_i=e_j\omega^j{}_i+n_\alpha\beta^\alpha{}_i$。这与第10章约定相同。

\begin{figure}[htbp]
\centering\DGDiagram{adapted-frame-block}

\caption{适配标架中，环境联络的对角块在切丛、法丛内求导，非对角块通过第二基本形式交换两类分量。上方块为$-\beta^{\mathsf T}$，负号由环境内积相容决定；此图表示联络分块，不将联络误当成逐纤维的线性映射。}
\label{fig:adapted-frame-block}
\end{figure}

\begin{proposition}[结构方程的三块]\label{prop:adapted-structure-equations}
令$\bar\Omega$为环境曲率形式沿$f$的拉回，分块下标$T,N$表示切、法输出与输入。则
\begin{align*}
 \dd\theta+\omega\wedge\theta&=0,\qquad \beta\wedge\theta=0,\\
 \Omega&=\bar\Omega_{TT}+\beta^{\mathsf T}\wedge\beta,\\
 \dd\beta+\tau\wedge\beta+\beta\wedge\omega&=\bar\Omega_{NT},\\
 \Omega^\perp&=\bar\Omega_{NN}+\beta\wedge\beta^{\mathsf T},
 \qquad\Omega^\perp=\dd\tau+\tau\wedge\tau.
\end{align*}
矩阵值形式的乘法按顺序进行，转置不交换楔积因子的先后。
\end{proposition}
\begin{proof}
在切标架中，无挠给第一条切结构方程。由$h^\alpha_{ij}=h^\alpha_{ji}$，
$\beta^\alpha{}_i\wedge\theta^i=h^\alpha_{ij}\theta^j\wedge\theta^i=0$。再将分块矩阵代入第10章已经证明的
$\bar\Omega=\dd\bar\omega+\bar\omega\wedge\bar\omega$：左上块为$\Omega-\beta^{\mathsf T}\wedge\beta$，左下块为$\dd\beta+\tau\wedge\beta+\beta\wedge\omega$，右下块为$\Omega^\perp-\beta\wedge\beta^{\mathsf T}$。移项即得。

例如切块在$(X,Y)$上为
\[
 (\beta^{\mathsf T}\wedge\beta)^i{}_j(X,Y)
 =\sum_\alpha\bigl(\beta^\alpha{}_i(X)\beta^\alpha{}_j(Y)
                 -\beta^\alpha{}_i(Y)\beta^\alpha{}_j(X)\bigr),
\]
即Gauss式的两项。交叉块包含$h$的导数和切、法联络修正，给Codazzi；法块给Ricci式。故这不是三组新增假设，而是同一环境结构方程在分解后的表达。
\end{proof}

\subsection*{由余标架求联络}
给出正交余标架$\theta$后，可以先解$\dd\theta+\omega\wedge\theta=0$和$\omega^{\mathsf T}=-\omega$，再求曲率，而不必先计算坐标Christoffel符号。存在性由Levi--Civita联络保证；唯一性也可直接看出。若两个解之差为$\eta^i{}_j=a_{ijk}\theta^k$，斜对称给$a_{ijk}=-a_{jik}$，零楔积给$a_{ijk}=a_{ikj}$。因此
\[
 a_{ijk}=-a_{jik}=-a_{jki}=a_{kji}=a_{kij}=-a_{ikj}=-a_{ijk},
\]
所以$\eta=0$。这正是活动标架计算中“解无挠及度量相容条件”的代数依据；本章习题用球面完整演示其符号与混合块。

\subsection*{相容条件为何足以在局部重建浸入}
下面证明欧氏环境的反向结论。它只需第10章平坦联络的局部平行标架，以及第7章闭一形式的局部恰当性，不预先调用子流形存在定理。

\Needspace{19\baselineskip}
\begin{theorem}[欧氏子流形基本定理的局部形式]\label{thm:local-immersion-reconstruction}
设$(M^m,g)$为黎曼流形，$E\to M$为带正定内积和相容联络$\nabla^E$的秩$r$向量丛，$B$为对称的$E$值二次协变张量。定义$A_\xi$由$g(A_\xi X,Y)=\langle B(X,Y),\xi\rangle$。若Gauss、Codazzi、Ricci三式对这些数据成立，且环境曲率项均取零，则每点有邻域$U$及等距浸入$F:U\to\R^{m+r}$，其法丛、法联络及第二基本形式在某个丛等距同构$E|_U\simeq NF$下恰为给定数据。缩小到连通邻域后，带该法丛识别的实现只差一个欧氏刚体运动。
\end{theorem}
\begin{proof}
在$\mathcal E=TM\oplus E$上定义
\[
 \mathcal D_X(Y,\xi)=
 \bigl(\nabla_XY-A_\xi X,\ B(X,Y)+\nabla_X^E\xi\bigr).
\]
方向线性及Leibniz规则由各项得到。对内积求导，混合项因$g(A_\xi X,Y)=\langle B(X,Y),\xi\rangle$成对抵消，所以$\mathcal D$相容。

计算$\mathcal D$的曲率：作用在$(Z,0)$上的切、法两块，恰为前节证明中的Gauss和Codazzi左端减去右端；作用在$(0,\xi)$上的法块为Ricci式的差。三组条件使这三块为零；剩下的切法交叉块是另一交叉块的负伴随，因为相容联络的曲率是斜伴随算子，故亦零。于是$\mathcal D$平坦。

由定理\ref{thm:flat-parallel-frame}，缩小到坐标球$U$后取$\mathcal D$平行正交归一标架。将其对应到$\R^{m+r}$的标准基，得到光滑丛等距映射$\Phi:\mathcal E|_U\to U\times\R^{m+r}$，满足
\[
 X(\Phi(s))=\Phi(\mathcal D_Xs).
\]
定义$\R^{m+r}$值一形式$\alpha(X)=\Phi(X,0)$。逐分量计算
\begin{align*}
 \dd\alpha(X,Y)
 &=\Phi\bigl(\mathcal D_X(Y,0)-\mathcal D_Y(X,0)-([X,Y],0)\bigr)\\
 &=\Phi\bigl(\nabla_XY-\nabla_YX-[X,Y],\ B(X,Y)-B(Y,X)\bigr)=0.
\end{align*}
把$U$取为坐标星形域，第7章Poincare引理逐分量给$F$使$\dd F=\alpha$。于是
\[
 \langle\dd F(X),\dd F(Y)\rangle=g(X,Y).
\]
正定性保证微分单射，$F$为等距浸入。$\Phi(0,E)$正好是$\dd F(TM)$的正交补；对$\dd F(Y)=\Phi(Y,0)$及$\Phi(0,\xi)$求导，$\mathcal D$定义正给所需的Gauss和Weingarten分解，因此所有法丛数据都被实现。

若另有实现，其环境平凡丛给出另一个保持$\mathcal D$的正交平凡化$\widetilde\Phi$。两平凡化之比在连通$U$上为常正交矩阵$Q$：对任意平行截面，两边系数均恒定。故$\dd\widetilde F=Q\dd F$，从而$\widetilde F=QF+b$，$b$为常向量。这包括反射，不只限于保定向刚体运动。
\end{proof}

特别地，对超曲面取$E=U\times\R$、平凡法联络及$B(X,Y)=g(SX,Y)n$，Ricci条件自动成立；Gauss和Codazzi已足以局部实现。结论是局部浸入，不能由此推出全局嵌入：平坦联络在非单连通区域可能有绕行，闭一形式也可能不恰当，且局部浸入延拓后可能自交。


\Needspace{17\baselineskip}
\section{极小子流形及基础例子}
\subsection*{体积元的微分}
本节对不定向流形也使用体积：在坐标中取$\sqrt{\det g}\,|\dd x^1\cdots\dd x^m|$，换卡时按Jacobian行列式的绝对值变换，再用单位分解积分。这一体积密度无需定向；在定向卡上它与第8章的黎曼体积形式给出相同积分。这里只建立紧区域体积所需的密度表达，不发展一般密度丛理论。
取光滑变分$f_t:M\to\bar M$，$f_0=f$，在紧区域上各$f_t$均为浸入。变分场$V=\partial_tf_t|_0$沿$f$取值，写$V=V^\top+V^\perp$；将切向部分经$\dd f$反识别为$M$的向量场。局部坐标下$g_{ij}(t)=\bar g(\partial_if_t,\partial_jf_t)$。环境无挠和度量相容给
\[
 \dot g_{ij}=\bar g(\bar\nabla_{\partial_i}V,\partial_jf)
             +\bar g(\partial_if,\bar\nabla_{\partial_j}V).
\]
Jacobi行列式公式给$\partial_t\sqrt{\det g_t}|_0=\tfrac12\operatorname{tr}(g^{-1}\dot g)\sqrt{\det g}$。在一点取正交归一基，便得
\[
 \left.\partial_t\dd\mathrm{vol}_{g_t}\right|_0
 =\left(\operatorname{div}_M V^\top-m\bar g(\mathbf H,V^\perp)\right)
   \dd\mathrm{vol}_g.
\]
这里散度项来自$\sum_i g(\nabla_{e_i}V^\top,e_i)$，法项来自
$\sum_i\bar g(\bar\nabla_{e_i}V^\perp,e_i)=-\sum_i\bar g(V^\perp,B(e_i,e_i))$。这一推导也固定了平均曲率的符号。

\Needspace{19\baselineskip}
\begin{theorem}[体积第一变分与极小性]\label{thm:first-volume-variation}
设$\Omega\Subset M$为具有光滑边界的紧区域，$\nu$为$\partial\Omega$在$M$内的外单位余法向，则
\[
 \left.\frac{\dd}{\dd t}\operatorname{Vol}_{f_t^*\bar g}(\Omega)\right|_0
 =-m\int_\Omega\bar g(\mathbf H,V^\perp)\,\dd\mathrm{vol}_g
   +\int_{\partial\Omega}g(V^\top,\nu)\,\dd\mathrm{vol}_{\partial\Omega}.
\]
故$\mathbf H=0$当且仅当对所有具有紧支变分场、固定区域边界的局部变分，体积一阶导数为零。这里极小表示驻定，不直接表示达到体积最小值。
\end{theorem}
\begin{proof}
在紧区域上可对积分号内求导，再对上面的体积元公式积分。若$M$定向，散度定理由第8章Stokes应用于$\iota_{V^\top}\dd\mathrm{vol}_g$得到；一般情形用局部定向图及单位分解，交叠时体积密度与外余法向通量同时相容，内部边界项相消，得到同一公式。因此不必给整个子流形选定向。

若$\mathbf H=0$且$V$在边界为零，两项均零。反向，在任意小坐标域内取非负紧支函数$\chi$，令$V=\chi\mathbf H$。这种沿$f$的场可实现为$f_t(p)=\exp_{f(p)}^{\bar M}(tV_p)$；在支集外保持$f$，紧性和浸入微分的开条件保证共同的小时间内仍为浸入。第一变分为$-m\int\chi|\mathbf H|^2$。若所有这种变分都驻定，连续性给$\mathbf H=0$。该论证只用小时间指数映射，不要求环境完备。
\end{proof}

对$m=1$，$\mathbf H=B(T,T)$，极小曲线正好是内蕴单位速曲线在环境中零加速度的情形，即环境测地线。对高维，零迹只约束$B$的总和，不要求各方向的弯曲为零。
\begin{figure}[htbp]
\centering\DGDiagram{mean-curvature-balance}

\caption{二维超曲面同一法空间中两个主方向的法向弯曲向量相加。左图$k_1=k_2=-1/a$给$\mathbf H=-n/a$；右图$k_1=1,k_2=-1$给该点$\mathbf H=0$而$B\ne0$。右图仅作逐点代数示意，不由一个点的零迹断言整张曲面极小。}
\label{fig:mean-curvature-balance}
\end{figure}

\subsection*{欧氏坐标函数与紧性障碍}
\begin{proposition}[浸入向量的Laplace方程]\label{prop:laplace-immersion}
对欧氏浸入$F:M^m\to\R^N$，按第11章$\Delta=\operatorname{div}\operatorname{grad}$的约定逐分量有
\[
 \Delta F=m\mathbf H,\qquad
 \Delta|F|^2=2m+2m\langle F,\mathbf H\rangle.
\]
因此$F$极小当且仅当全部欧氏坐标函数调和。不存在正维、紧致、无边界的欧氏极小浸入。
\end{proposition}
\begin{proof}
函数Hessian的迹在局部正交归一标架中为
\[
 \Delta F=\sum_i\bigl(D_{e_i}(\dd F(e_i))-\dd F(\nabla_{e_i}e_i)\bigr)
          =\sum_i B(e_i,e_i).
\]
逐分量的乘法法则给$\Delta|F|^2=2\sum_i|\dd F(e_i)|^2+2\langle F,\Delta F\rangle$；等距性使平方和为$m$。若紧致无边界且$\mathbf H=0$，对第二式积分，散度定理左端为零，右端为$2m\operatorname{Vol}(M)>0$，矛盾。不定向情形仍用体积密度版本；有边界时会产生边界通量，不能沿用矛盾。
\end{proof}

球面嵌入的$\mathbf H=-F/a^2$，所以$\Delta F=-mF/a^2$，与第11章Laplace符号一致。欧氏极小曲面有$\mathrm{Scal}=-|B|^2$；二维时$2K=-|B|^2$，故非全测地处$K<0$。

\subsection*{极小图像方程及一个有限范围的最小性结论}
设$u:U\subset\R^m\to\R$光滑，取图像$F(x)=(x,u(x))$及向上法向$n=(-Du,1)/W$，$W=(1+|Du|^2)^{1/2}$。直接计算
\[
 g_{ij}=\delta_{ij}+u_i u_j,\quad
 g^{ij}=\delta_{ij}-\frac{u_i u_j}{W^2},\quad
 II_{ij}=\frac{u_{ij}}{W},\quad
 mH=\frac1W\left(\delta_{ij}-\frac{u_i u_j}{W^2}\right)u_{ij}.
\]
逆矩阵可乘回$g$验证；$II_{ij}$由$\langle\partial_{ij}F,n\rangle$给出。展开$\partial_i(u_i/W)$，恰有
\[
 mH=\operatorname{div}_{\R^m}\left(\frac{Du}{\sqrt{1+|Du|^2}}\right).
\]
所以极小图像方程为右端等于零。$m=2$时等价于
\[
 (1+u_y^2)u_{xx}-2u_xu_yu_{xy}+(1+u_x^2)u_{yy}=0.
\]
例如$u=xy$在原点有$H=0$，却在一般点给分子$-2xy$，故不为极小图像。

设$\Omega\Subset U$边界光滑，$u$在其闭包光滑并满足此方程。此时它确实在具有同一边值的光滑图像中使面积最小：函数$L(p)=\sqrt{1+|p|^2}$的Hessian为
\[
 D^2L(p)=\frac{I}{\sqrt{1+|p|^2}}-
 \frac{pp^{\mathsf T}}{(1+|p|^2)^{3/2}},
\]
对非零$q$有$q^{\mathsf T}D^2L(p)q\geq |q|^2/(1+|p|^2)^{3/2}>0$。沿$p$到$q$的线段积分得到凸性不等式$L(q)\geq L(p)+DL(p)\cdot(q-p)$。令$p=Du,q=Dv$并积分，若$v=u$在边界，则
\[
 \operatorname{Area}(v)-\operatorname{Area}(u)
 \geq\int_\Omega\frac{Du}{W}\cdot D(v-u)=0,
\]
末步分部积分用了极小方程和零边值。这是固定图像类中的最小性，不据此宣称对所有浸入或所有拓扑类型的竞争者都有全局最小性。一般极小子流形的第二变分、稳定性和Plateau问题留待另行建立。


\section{子流形几何习题及详细解答}
13.1将梁书4.4注3的球面投影练习改述并扩展（印刷101/PDF117）；13.2将定理4-4-2的正则水平集主题练习化（印刷99/PDF115）。两题均为文内主题改编，不冒充章末题号。13.3--13.8为原创补充。已查看梁书第4章原题17（印刷104/PDF120），它讨论Lorentz类光超曲面，留到本书第21章的相应范围，不套入本章正定分解。

\Needspace{12\baselineskip}
\subsection*{习题13.1：球面投影与诱导度量的两个定义域}
梁书4.4注3改述并扩展。在$\R^3\setminus\{0\}$球坐标中写出欧氏度量；对$S^2(a)$构造环境切空间上的正交投影$P$及退化双线性形式$h(v,w)=\langle Pv,Pw\rangle$。说明它与球面黎曼度量的关系，计算球面的$S,\mathbf H,K$。
\begin{proof}[解答]
欧氏度量为$\dd r^2+r^2\dd\theta^2+r^2\sin^2\theta\,\dd\phi^2$。球面外单位法向$n=\partial_r$，法余矢量$n^\flat=\dd r$，故
\[
 Pv=v-\langle v,n\rangle n,\quad
 P^2=P,\quad\ker P=\R n,\quad\operatorname{im}P=TS^2(a),
 \qquad h=\bar g-\dd r\otimes\dd r.
\]
$h$在环境三维切空间上有一维核；沿包含映射拉回后则为$a^2(\dd\theta^2+\sin^2\theta\,\dd\phi^2)$，在球面二维切空间上正定。不能把环境退化说成诱导球面度量退化。

在笛卡儿坐标$n=F/a$，$\dd n(X)=X/a$，所以$S=-I/a$、$B=-gn/a$、$\mathbf H=-n/a$，Gauss式给$K=a^{-2}$。坐标极点改用其他卡；投影公式本身无极点奇性。换法向只改变标量$II,H$和$S$的符号，$B,\mathbf H,K$不变。
\end{proof}

\Needspace{12\baselineskip}
\subsection*{习题13.2：正则水平集的Hessian公式}
梁书定理4-4-2主题改编并扩展。设$(\bar M,\bar g)$上$\varphi$光滑，$M=\varphi^{-1}(c)$上$\operatorname{grad}\varphi\ne0$。取$n=\operatorname{grad}\varphi/|\operatorname{grad}\varphi|$，用环境Hessian表示$II$。若用$\psi=\rho(\varphi)$代替$\varphi$且$\rho'(c)\ne0$，几何量如何改变？
\begin{proof}[解答]
对切向$Y$有$Y\varphi=0$，故$\operatorname{grad}\varphi$正交于$TM$。记$q=|\operatorname{grad}\varphi|$，对切向$X,Y$有
\[
 \bar g(\bar\nabla_Xn,Y)=q^{-1}\operatorname{Hess}_{\bar g}\varphi(X,Y),
 \qquad II(X,Y)=-q^{-1}\operatorname{Hess}_{\bar g}\varphi(X,Y).
\]
求导$q^{-1}$的项平行于$\operatorname{grad}\varphi$，与$Y$正交，故没有遗漏。亦可对$Y\varphi=0$求$X$导数得到同式。

链式法则给$\operatorname{Hess}\psi=\rho'\operatorname{Hess}\varphi+\rho''\dd\varphi\otimes\dd\varphi$，最后一项在切向槽为零。$|\operatorname{grad}\psi|=|\rho'|q$，于是$n_\psi=\operatorname{sgn}(\rho')n$，$II_\psi=\operatorname{sgn}(\rho')II$，但$B=II\,n$不变。正则性仅保证附近同一水平集分支；若$\rho$在远处不单射，$\psi^{-1}(\rho(c))$可能另有分支。取$\varphi=|x|^2$在球面验证$II=-g/a$，排除整体符号错误。
\end{proof}

\Needspace{12\baselineskip}
\subsection*{习题13.3：圆柱的内蕴平坦与外在弯曲}
原创。对$F(\theta,z)=(a\cos\theta,a\sin\theta,z)$计算度量、第二基本形式、主曲率及平均曲率向量。哪些圆柱测地线也是环境直线？
\begin{proof}[解答]
$g=a^2\dd\theta^2+\dd z^2$。取外法向$n=(\cos\theta,\sin\theta,0)$及正交单位基$e_1=a^{-1}\partial_\theta,e_2=\partial_z$，则
\begin{align*}
 S&=\operatorname{diag}(-a^{-1},0),\qquad B(e_1,e_1)=-n/a,\\
 B(e_1,e_2)&=B(e_2,e_2)=0,\qquad\mathbf H=-n/(2a),\qquad K=0.
\end{align*}
因此平坦不意味着全测地或极小。内蕴仿射测地线在局部提升角度后为$\theta=At+B,z=Ct+D$。环境加速度为$-aA^2n$，所以成为环境测地线当且仅当$A=0$，即直母线（含常值曲线）。圆周$C=0,A\ne0$虽是圆柱测地线，却不为欧氏直线；一般$A,C\ne0$给螺线。
\end{proof}

\Needspace{12\baselineskip}
\subsection*{习题13.4：余维二的非平坦法丛}
原创。令$F(x,y)=(x,y,(x^2-y^2)/2,xy)$。证明它是欧氏极小浸入，且在原点法丛曲率非零。用Ricci方程核对其符号。
\begin{proof}[解答]
设$q=1+x^2+y^2$，直接内积得$g=q(\dd x^2+\dd y^2)$。二维共形度量的Laplace为$q^{-1}(\partial_{xx}+\partial_{yy})$；四个分量均为欧氏调和函数，故$\Delta_gF=0$，由命题\ref{prop:laplace-immersion}，$\mathbf H=0$。$q>0$保证全域浸入。

原点切基为标准$e_1,e_2$，法基为$n_1=e_3,n_2=e_4$；二阶导数的法分量给
\[
 A_{n_1}=\begin{pmatrix}1&0\\0&-1\end{pmatrix},\qquad
 A_{n_2}=\begin{pmatrix}0&1\\1&0\end{pmatrix},\qquad
 [A_{n_1},A_{n_2}]=\begin{pmatrix}0&2\\-2&0\end{pmatrix}.
\]
故$\langle R^\perp(e_1,e_2)n_1,n_2\rangle=-2$，法联络不平坦；两个形算子不能同时对角化。Gauss式给$K=\langle n_1,-n_1\rangle-|n_2|^2=-2$，$|B|^2=4$，恰有$2K=-|B|^2$。这是真正的极小曲面，不能从极小性推断法丛平坦。
\end{proof}

\Needspace{12\baselineskip}
\subsection*{习题13.5：Clifford环面的环境依赖}
原创。令$F(u,v)=2^{-1/2}(\cos u,\sin u,\cos v,\sin v)$，两参数均模$2\pi$。证明它在单位$S^3$中极小、内蕴平坦，在$\R^4$中却不极小，并说明紧性障碍为何不矛盾。
\begin{proof}[解答]
$|F|=1$，$g=(\dd u^2+\dd v^2)/2$。在$S^3$内的单位法向为
\[
 n=2^{-1/2}(\cos u,\sin u,-\cos v,-\sin v).
\]
它与$F,F_u,F_v$正交。$n_u=F_u$、$n_v=-F_v$，且这些导数切于$S^3$，故$S$的两个特征值为$-1,1$，平均值零，$F$在$S^3$内极小。环境截面曲率为1，Gauss式给$K=1-1=0$，与常系数度量一致。

在$\R^4$中，单位球面有$B^{S^3}(X,Y)=-g(X,Y)F$。连续应用两层Gauss分解，环面的欧氏第二基本形式为$B^{\R^4}=B^{S^3\text{内}}-gF$，取一半迹得$\mathbf H_{\R^4}=-F\ne0$。也可从$\Delta_gF=2(F_{uu}+F_{vv})=-2F$核对。紧致无边界欧氏极小浸入不存在的命题要求欧氏环境，此例的极小环境是弯曲的$S^3$。
\end{proof}

\Needspace{12\baselineskip}
\subsection*{习题13.6：螺旋面不是平面却处处极小}
原创。对$c\ne0$，令$F(u,v)=(u\cos v,u\sin v,cv)$。计算$g,II,H,K$，并验证$u=0$不导致参数奇性。
\begin{proof}[解答]
取$W=\sqrt{u^2+c^2}$，有$g=\dd u^2+W^2\dd v^2$，
$n=(c\sin v,-c\cos v,u)/W$。直接求二阶导数得
\[
 II_{uu}=II_{vv}=0,\qquad II_{uv}=II_{vu}=-c/W.
\]
因此$\operatorname{tr}_gII=0$，处处极小；而
\[
 K=\frac{\det II}{\det g}=-\frac{c^2}{(u^2+c^2)^2}<0.
\]
相对正交基$\partial_u,W^{-1}\partial_v$，$S$非对角项为$-c/W^2$，特征值为$\pm c/W^2$（作为无序对）；二者互相抵消但不为零，所以不全测地。由于$W\geq|c|>0$，在$u=0$时$F_v=(0,0,c)$仍非零且与$F_u$独立。若令$c=0$，这一结论失效，不能将退化参数当同一个正则例子。
\end{proof}

\Needspace{12\baselineskip}
\subsection*{习题13.7：球面活动标架的符号核对}
原创。在$S^2(a)$非极点坐标卡取$\theta^1=a\dd\theta$、$\theta^2=a\sin\theta\dd\phi$，求切联络矩阵、曲率矩阵和外法向混合块，逐项验证环境欧氏曲率为零。
\begin{proof}[解答]
由$\dd\theta^1=0$、$\dd\theta^2=(\cot\theta/a)\theta^1\wedge\theta^2$及无挠、斜对称条件，得到
\[
 \omega^1{}_2=-\cos\theta\,\dd\phi,\qquad
 \omega^2{}_1=\cos\theta\,\dd\phi,\qquad
 \Omega^1{}_2=\sin\theta\,\dd\theta\wedge\dd\phi
             =a^{-2}\theta^1\wedge\theta^2.
\]
其余由斜对称决定，故$\Omega^1{}_2(e_1,e_2)=\langle R(e_1,e_2)e_2,e_1\rangle=K=a^{-2}$。这里列指标为被求导的向量，不与另一种常见行指标约定混用。

外法向使$\beta=(-\theta^1/a,-\theta^2/a)$，$\tau=0$。于是$(\beta^{\mathsf T}\wedge\beta)^1{}_2=a^{-2}\theta^1\wedge\theta^2=\Omega^1{}_2$，切块$\bar\Omega_{TT}=0$；$\dd\beta+\beta\wedge\omega=0$由无挠结构方程成立；法块$\beta\wedge\beta^{\mathsf T}=0$。因此全部环境曲率块为零。曲率形式的符号若与本题相反，应先检查指标次序和$R$定义，而非改变球面正截面曲率。
\end{proof}

\Needspace{12\baselineskip}
\subsection*{习题13.8：局部相容不等于全局单值重建}
原创。在$M=\R/(2\pi\mathbb Z)$上给$g=\dd s^2$，令$E$为平凡单位法线丛、法联络平凡且$B=0$。验证三组欧氏相容方程，证明局部能浸入$\R^2$，但不存在全局实现这些数据的浸入。指出基本定理证明中的障碍。
\begin{proof}[解答]
一维切丛曲率为零，平凡法联络亦平坦，$B,A,\nabla B$全部为零，所以三组条件均成立。小弧上的$F(s)=(s,0)$是局部等距实现。

若全局实现存在，取覆盖参数$s\in\R$，因$B=0$且内蕴$\partial_s$平行，Gauss式给$F''(s)=0$，所以$F(s)=p+sv$，$|v|=1$。但周期性要求$F(s+2\pi)=F(s)$，即$2\pi v=0$，矛盾。

此例$TM\oplus E$的平坦联络已有全局平行标架，障碍不是法丛绕行，而是重建一形式$\alpha=(\dd s,0)$的周期$\int_{S^1}\alpha=(2\pi,0)$非零，闭而不恰当。因此局部相容不能省去积分的全局条件；连通或局部唯一性也不能代替它。
\end{proof}

